% Generated draft: requires mathematical and source review.
\documentclass[11pt]{article}
\usepackage[margin=25mm]{geometry}
\usepackage{fontspec}
\setmainfont{Latin Modern Roman}
\usepackage{amsmath,amssymb,mathtools,unicode-math}
\setmathfont{Latin Modern Math}
\usepackage{enumitem,needspace,xurl,hyperref}
\setlength{\parindent}{0pt}
\setlength{\parskip}{5pt}
\setlength{\emergencystretch}{3em}
\newcommand{\N}{\mathbb{N}}
\newcommand{\Z}{\mathbb{Z}}
\newcommand{\Q}{\mathbb{Q}}
\newcommand{\R}{\mathbb{R}}
\newcommand{\C}{\mathbb{C}}
\newcommand{\F}{\mathbb{F}}
\newcommand{\A}{\mathbb{A}}
\newcommand{\PP}{\mathbb{P}}
\newcommand{\E}{\mathbb{E}}
\newcommand{\eps}{\varepsilon}
\newcommand{\dd}{\,\mathrm{d}}
\begin{document}
\section*{UnsolvedMath ID: 20000243}
\textbf{Source problem code:} catalogue entry
\subsection{Definitions and mathematical statement}
The primary question is attributed to Anastasia Stavrova. For a base scheme $S$, let $\mathrm{Sm}_S$ denote smooth $S$-schemes. The usual construction of unstable motivic homotopy theory uses simplicial presheaves on $\mathrm{Sm}_S$, Nisnevich descent, and the identifications $X\times_S\A^1_S\simeq X$; a stable version also requires specifying a stabilization, usually by the pointed projective line. These constructions clarify the meaning of extension, but do not specify which theorems from the field case should persist.

\textbf{Original question.} Can $\A^1$-homotopy theory be extended over general base schemes, in particular over a discrete valuation ring? Explain the issues if the desired extension fails. The extracted phrase ``general smooth schemes'' leaves the intended base class and the scheme relative to which smoothness is required unspecified. The original remark, attributed to Aravind Asok, warns that many basic results fail, specifically the stable $\A^1$-connectivity theorem. Thus the question includes the behavior of field-level theorems, not merely the existence of categories.

\textbf{Attached research-aid claims, requiring independent verification.} The aid reports that unstable and stable motivic categories can be constructed over arbitrary schemes, whereas published stable connectivity estimates over a one-dimensional base permit a loss of one degree. It also mentions an analogous dimension-loss estimate in an unidentified 2025 unstable preprint. The notions of connectivity, precise hypotheses, indexing, and references for these estimates are not supplied; these assertions cannot be turned here into a fully quantified connectivity theorem.

The aid gives the following more specific proposed reduction. Let $k$ be an infinite perfect field, $R=(k[t]_{(t)})^h$ its indicated henselization, $K=\operatorname{Frac}(R)$, and identify its residue field with $k$. Write $K_q^M$ for Milnor $K$-theory, with $K_0^M(k)=\Z$. Let $i_q:K_q^M(R)\to K_q^M(K)$ and $r_q:K_q^M(R)\to K_q^M(k)$ be induced by inclusion and reduction, and choose the coefficient-field section $s_q:K_q^M(k)\to K_q^M(K)$. Normalize the tame residue by
\[\partial_t\bigl(\{t\}\cdot s_{q-1}(\gamma)\bigr)=\gamma.\]
For every integer $q\ge1$, the report asserts that $i_q$ is injective with image $\ker\partial_t$, so that $u_q(\alpha)\in K_q^M(R)$ is uniquely defined by
\[i_q(u_q(\alpha))=\alpha-\{t\}\cdot s_{q-1}(\partial_t\alpha).\]
Its claimed split exact sequence is
\[0\longrightarrow K_q^M(R)\xrightarrow{(i_q,r_q)}K_q^M(K)\oplus K_q^M(k)\xrightarrow{\Phi_q}K_{q-1}^M(k)\oplus K_q^M(k)\longrightarrow0,\]
\[\Phi_q(\alpha,\beta)=\bigl(\partial_t\alpha,\,\beta-r_q(u_q(\alpha))\bigr).\]
A splitting described by these choices sends $(\gamma,\delta)$ to $(\{t\}\cdot s_{q-1}(\gamma),\delta)$. The two reported obstruction components are the shifted tame residue and the generic/special mismatch in degree $q$. This is an unreviewed, choice-dependent Milnor $K$-theory claim in this particular equicharacteristic henselian example; the source supplies no implication making it a solution of the general motivic extension question.
\subsection{Short English statement}
Stavrova asks how $\A^1$-homotopy theory extends beyond fields, especially to a DVR; Asok highlights stable connectivity as an obstacle. The attached unreviewed aid distinguishes construction of motivic categories from connectivity estimates and proposes a split Milnor $K$-theory obstruction sequence for one equicharacteristic henselian DVR.
\subsection{Sources}
\begin{itemize}
\item Local UnsolvedMath registry, permanent ID 20000243, source code catalogue entry; canonical location \texttt{aim-algebraic-geometry-notes.json}, notes[242], tag \texttt{question}.
\item AIM workshop, \emph{Projective modules and $\A^1$-homotopy theory}, Question 6, attributed to Anastasia Stavrova; accompanying remark attributed to Aravind Asok. Supplied source URL: \url{https://aimath.org/pastworkshops/projectiveA1problems.pdf}. No PDF was consulted for this draft.
\item Attached research attempt 1: result status \texttt{partial\_result}, problem status \texttt{partially\_solved}, rigor label \texttt{proved}, novelty confidence low. The source explicitly calls it a machine-generated, non-peer-reviewed research aid requiring independent verification before citation. Its references to published stable bounds and a 2025 unstable preprint contain no bibliographic identification; the original literature field is empty.
\end{itemize}
\end{document}
